
\documentclass[11pt]{article}
\usepackage[T1]{fontenc}
\usepackage[utf8]{inputenc}
\usepackage[english]{babel}
\usepackage{amsmath,amssymb,mathtools,array,graphicx}
\usepackage{geometry}
\usepackage{microtype}
\geometry{a4paper,margin=2.55cm}
\setlength{\parindent}{1.25em}
\setlength{\parskip}{0pt}
\setlength{\emergencystretch}{8em}
\allowdisplaybreaks
\raggedbottom
\graphicspath{{figures/}}
\newcommand{\A}{\Delta}
\newcommand{\sourceplate}[1]{\begin{center}\includegraphics[width=0.985\linewidth]{#1}\end{center}}
\newcommand{\widebar}[1]{\overline{#1}}
\begin{document}

The relatively complete system modulo $\nu$ therefore consists of the six forms:
\[
f,\ \theta,\ K,\ j,\ \A,\ N.
\]

As an example of the sequence expansion in \S\ 2 we give the derivation of formula (3.)a in full; in later calculations the final formula III will be set down directly. (Polars of vanishing forms have already been omitted during the calculation; the first line gives the values inserted according to Expansion II, and the second line gives the final values found recursively, beginning with the last equation of the system.) It follows, according to Expansion II:

\sourceplate{p50_derivation_formulas_1_2.png}
\sourceplate{p51_derivation_formulas_3_9.png}
\sourceplate{p52_derivation_formulas_and_note.png}

\[
\tag{4}
g_\nu^3=-\frac7{10}\theta_\nu^3+\frac12u_x^2\theta,
\qquad
g_\nu^4=\frac35\theta.
\]

From (3.) and (4.) it follows, by transfer from the binary domain, that
\[
(\theta\theta'u)^2=\frac13\,j\cdot f\pmod{\nu}.\footnotemark
\]
The remaining forms of the form sequence $(\theta\theta'u)^2$ and the form $\theta'_\nu$ are reducible to symbols $\nu$. We can therefore replace:
\[
\bmod\ \nu:\quad \text{foldings with } f\cdot j
\text{ by the corresponding foldings with }(\theta\theta'u)^2.
\]
Furthermore,
\[
(\vartheta\vartheta'x)^2=\frac43 f\cdot\A,
\qquad
\theta_{g\nu}(\vartheta\vartheta'x)=-N.
\]
\footnotetext{Gordan--Kerschensteiner, loc. cit., p. 181.}

\begin{center}
{\Large\bfseries Chapter III.\quad The relatively complete system modulo $(s,\varrho,t)$.}
\end{center}

\begin{center}
\textbf{\S\ 7.\quad Overview of the formation of the system.}
\end{center}

In order to pass from the relatively complete system modulo $\nu$ to the relatively complete system with the next higher modulus, one has, by Definition \S\ 3, to form the system of $\nu$ with respect to this higher modulus and to fold it with the forms of the relatively complete system modulo $\nu$. But $\nu=u_\nu^4$, as a contravariant of the fourth degree, is dual to the form $f$. If therefore we regard $\nu$ as the ground form, we obtain a relatively complete system of six forms modulo $\nu(\nu)=(\nu\nu_1x)^4=s_x^4$.

Thus, for the formation of the relatively complete system modulo $s$ (System III), we have to transvect
\[
\begin{array}{ll}
\text{System I:} & f,\ \theta,\ K,\ j,\ \A,\ N\quad \text{over}\vphantom{\dfrac11}\\[0.3em]
\text{System II:} & \nu,\ H=(\nu\nu_1x)^2u_\nu^2u_{\nu_1}^2,\ L=(\nu\eta x)H_x^2u_\nu^3u_\eta^3,\\
& g=(\nu\eta x)^4H_x^2,\ \sigma=H_\eta^2u_\nu^2u_\eta^4,\ (\nu\sigma x).
\end{array}
\]
It will be shown (formula (13.)) that the form $g$ is reducible, and hence that System II consists of only five forms.

In the course of the calculation the quadratic forms
\[
\varrho=u_\varrho^2=\theta_\nu^4u_\vartheta^2,
\qquad
 t=t_x^2=a_\eta^4H_x^2
\]
are adjoined as moduli, so that one obtains a relatively complete system modulo $(s,\varrho,t)$; and the modulus $(\varrho,t)$ is to be regarded as a higher modulus than the modulus $s$.

The arrangement of the individual forms is to follow the order in the coefficients. We normalize the folding
\[
\theta_\nu(K_\nu,\theta_\nu,\ldots)
\]
as higher than the folding
\[
H_y(H_y,L_y,\ldots).
\]
Consequently, by \S\ 1 and \S\ 3b, the order of the individual forms of the same order is uniquely determined.

The formation of the forms of the same order is carried out in tabular form:

I. Indication of which forms of Systems I and II are to be folded with each other in order to obtain the prescribed order in the coefficients.

II. Listing of the irreducible forms according to the scheme of the form sequence.

III. Reduction of the reducible or decomposable form sequences or forms, that is, of those places where the total form sequence breaks off, thereby proving that all irreducible constructions have been exhausted by those indicated.

IV. Consequences: 1) indication of newly arising reducers; 2) indication of those forms which do not enter, in products with forms of the same system, into folding with forms of the other system.

It will be shown that only the following occur:

1) foldings of one form of System I with one form of System II;

2) foldings of powers of $\theta$, respectively $H$, or of two forms of one system with one form of the second system.

We shall obtain the following scheme for folding:
\[
\begin{array}{r@{\ }c@{\ }l@{\qquad}c@{\qquad}r@{\ }c@{\ }l}
\multicolumn{3}{c}{\text{System I}} & \text{folded with} & \multicolumn{3}{c}{\text{System II}}\\[0.25em]
1. & \text{order} & f & & 2. & \text{order} & \nu\\
2. & ,, & \theta & & 4. & ,, & H\\
3. & ,, & K,j,\A & & 6. & ,, & L,\sigma\\
4. & ,, & N,\theta^2,f\cdot j & & 8. & ,, & (\nu\sigma x),H^2\\
5. & ,, & \theta\cdot K & & 10. & ,, & H\cdot L\\
6. & ,, & \theta^3,\A\cdot j & & 12. & ,, & H^3.
\end{array}
\]

For checking the calculation, besides the ``double reduction'', the two special forms serve:

\sourceplate{p55_special_forms_formulas_5_6.png}

\[
\tag{7}
\begin{aligned}
\text{a)}\quad (\nu\sigma x)^2
&=\frac12(Hsu)^2-\frac25\nu\cdot s_\nu^2-
\frac45u_x\,s_\eta(Hsu)^2
+u_x^2\left\{\frac16J\cdot\nu-\frac1{40}(ss'u)^4\right\},\\
\text{b)}\quad (\nu\sigma x)^3&=-\frac3{10}s_\eta(Hsu),\\
\text{c)}\quad (\nu\sigma x)^4&=\frac{21}{10}s_\eta^2-\frac3{10}Z-
\frac65u_xs_\eta^3+\frac1{10}u_x^2s_\eta^4.
\end{aligned}
\]
\[
\tag{8}
\begin{aligned}
\text{a)}\quad g_\nu&=-s_\eta+u_x\left(\frac92s_\eta^2-\frac12Z\right)-3u_x^2s_\eta^3+\frac12u_x^3s_\eta^4,\\
\text{b)}\quad g_\nu^2&=\frac{21}{10}s_\eta^2-\frac3{10}Z-2u_xs_\eta^3+\frac12u_x^2s_\eta^4,\\
\text{c)}\quad g_\nu^3&=-\frac7{10}s_\eta^3+\frac12u_xs_\eta^4,\\
\text{d)}\quad g_\nu^4&=\frac35s_\eta^4.
\end{aligned}
\]

\begin{center}
\textbf{\S\ 8.\quad Forms of the third order (System III).}
\end{center}

\noindent\textbf{Folding of $f$ with $\nu$.}

\noindent\emph{Irreducible forms:}
\[
\begin{array}{cccc}
a_\nu&\cdot&\cdot&\cdot\\
\cdot&a_\nu^2&\cdot&\cdot\\
\cdot&\cdot&\cdot&\cdot\\
\cdot&\cdot&\cdot&a_\nu^4=i.
\end{array}
\]

\noindent\emph{Reduction:}
\[
a_\nu^3=\frac13 i u_x\qquad\text{(formula (2.)).}
\]
\noindent\emph{Consequences:} 1) $a_\nu^3$ is a reducer. 2) $f$ enters into folding with $\nu$ only in products with contragredient forms $(j)$. The same holds for $\nu$ with respect to $f$.

\begin{center}
\textbf{\S\ 9.\quad Forms of the fourth order (System III).}
\end{center}

\noindent\textbf{Folding of $\theta$ with $\nu$.}

\noindent\emph{Irreducible forms:}
\[
\begin{array}{cccc}
(\theta\nu x)&\theta_\nu&\cdot&\cdot\\
(\theta\nu x)^2&\theta_\nu(\theta\nu x)&\theta_\nu^2&\cdot\\
\cdot&\theta_\nu(\theta\nu x)^2&\cdot&\theta_\nu^3.
\end{array}
\]

\noindent\emph{Reductions:} From the reducer $a_\nu^3$, by double reduction of the expressions $a_\nu^3(abu)$, $a_\nu^3b_\nu$, $a_\nu^3b_\nu(abu)$ according to the ansatz:

\sourceplate{p57_reduction_derivation_and_formula_9.png}

The formulas obtained are:
\[
\tag{9}
\begin{array}{c|c|c}
\theta_\nu^2(\theta\nu x)=0
& \theta_\nu^3(\theta\nu x)=0
& \theta_\nu^4u_x^2=u_\varrho^2=\varrho\quad(\text{adjoined modulus, \S\ 7}),\\[0.6em]
\theta_\nu^2(\theta\nu x)^2=\dfrac49 i f-\dfrac16s
& &
\end{array}
\]

\noindent\emph{Consequences:} 1) $\theta_\nu^2(\theta\nu x)^2$ is a reducer. 2) $\nu$ enters into folding with $\theta$ only in products with contragredient forms $(g)$.

\begin{center}
\textbf{\S\ 10.\quad Forms of the fifth order (System III).}
\end{center}

\[
\text{Folding of }\left\{\begin{array}{c}K,j,\A\text{ with }\nu,\\ f\text{ with }H.
\end{array}\right.
\]

\noindent\emph{Irreducible forms:}
\sourceplate{p58_irreducible_forms_order5.png}

\noindent\emph{Reductions:}

A) Form sequence $K_\nu^3$: reducer $a_\nu^3$.

Form $K_\nu^2(k\nu x)^2$: reducer $\theta_\nu^2(\theta\nu x)^2$.
\[
(k\nu x),\quad K_\nu(k\nu x),\quad K_\nu^2(k\nu x)
\]
are decomposable forms by \S\ 3.

B) Form
\[
\begin{aligned}
(j\nu x)^4&=(\widehat{a}\theta\nu x)^4
=a_\nu^4\theta+\theta_\nu^4f-4a_\nu^3\theta_\nu
-4a_\nu\{a_\nu\theta_x-(\widehat{a}\theta\nu x)\}^3\\
&\quad+6a_\nu^2\{a_\nu\theta_x-(\widehat{a}\theta\nu x)\}^2
=\varrho\cdot f+\frac53 i\theta-6a_\nu^2(\widehat{a}\theta\nu x)^2+4a_\nu(\widehat{a}\theta\nu x)^3,
\end{aligned}
\]
and therefore
\[
(j\nu x)^4=\varrho\cdot f+\frac53 i\theta\pmod{(\A,\nu)},
\qquad\text{(cf. the end of \S\ 5).}
\]

C) Form $\A_\nu^4$: reducer $\theta_\nu^4$.

Forms $\A_\nu^2$ and $\A_\nu^3$: reducer $a_\nu^3$ or $\theta_\nu^4$ by replacing the form $\A$ by lower forms of the form sequence (end of \S\ 5), that is, by double reduction of the expressions
\[
a_\vartheta a_\nu^3,
\qquad a_\vartheta a_\nu^3\theta_\nu,
\qquad a_\vartheta\theta_\nu^4.
\]
The double computation of $\A_\nu^3$ gives occasion for the reduction of the higher form $a_\eta^3$.

\noindent\emph{Reduction formulas:}
\[
\tag{10}
\begin{aligned}
\text{a)}\quad \A_\nu^2&=\frac35a_\eta^2-\frac3{10}(asu)^2+\frac13 i\theta+\frac15u_x^2(2a_\varrho^2-t),\\
\text{b)}\quad \A_\nu^3&=-\frac3{10}a_\varrho+\frac35u_x\left\{\frac{13}{10}a_\varrho^2-\frac25t\right\},\\
\text{c)}\quad \A_\nu^4&=a_\varrho^2-\frac25t.
\end{aligned}
\]

\noindent\emph{Derivation of the formulas:}
\sourceplate{p59_formula10_derivation.png}

D) Reduction of $a_\eta^3$ by double reduction of $\A_\nu^3$ (C).

The remaining reductions are by a calculation dualistically opposed to that of \S\ 9.

\noindent\emph{Reduction formulas:}
\[
\tag{11}
\begin{array}{rl|rl}
a_\eta^2(aHu)&=-\dfrac16(asu)^3,
& a_\eta^2(aHu)^2&=\dfrac49 i\nu-\dfrac16(asu)^4,\\[0.6em]
a_\eta^3&=-a_\varrho+\dfrac35u_xa_\varrho^2+\dfrac15u_xt,
& a_\eta^3(aHu)&=0,\\[0.6em]
& & a_\eta^4&=t\quad\text{(as adjoined modulus).}
\end{array}
\]

\noindent\emph{Consequences:}

1) $(j\nu x)^4$, $\A_\nu^2$, $a_\eta^2(aHu)$ are reducers.

2) $\nu$ enters only in products with contragredient forms into folding with $K,j,\A$; the same holds for $f$ with respect to $H$ and for $j$ with respect to $\nu$, since $\theta_\nu(\theta\cdot j,\nu,x^3)$ becomes reducible by \S\ 11 B.

\end{document}
